\documentclass[12pt]{article}
\usepackage[user]{cqd}

%% amsmath já carregado
%%\RequirePackage[utf8]{inputenc} já carregado
%\RequirePackage{etoolbox} já carregado
%\RequirePackage{color} já carregado
%\RequirePackage{calc} já carregado
%\RequirePackage[english]{babel} já carregado
%\RequirePackage{graphicx} já carregado
%\RequirePackage{geometry} já carregado

%%% espaço para definir os comandos \newtheorem, por exemplo:
%% %%  - theorem, corolary, lemma, proposition, definition, example etc
%%
%%%%%%%%%%%%%%%%%%%%%%%%%%%%%%%%%%%%
%% preencha as informações abaixo %%

%\titulo{Título do artigo}
\tituloingles{Título e artigo em inglês}

\autor{Nome Sobrenome 1}{Instituto/Unidade 1}{Instituição 1}{email@server.com.br}{Sobrenome 1, N.}
\autor{Nome Sobrenome 2}{Instituto/Unidade 2}{Instituição 2}{email@server.com.br}{Sobrenome 2, N.}
\autor{Nome Sobrenome 3}{Instituto/Unidade 3}{Instituição 3}{email@server.com.br}{Sobrenome 3, N.}
%\autor{Nome Sobrenome 4}{Instituto/Unidade 4}{Instituição 4}{email@server.com.br}{Sobrenome 4, N.}


%\resumo{Este resumo deve conter de 800 a 1100 caracteres, e deve enunciar de forma clara e sintética o problema de pesquisa, a abordagem metodológica empreendida, resultados e conclusões. \lipsum[1] }

%\palavraschave{palavras. separadas. por. ponto final.}


\abstract{O abstract do artigo deve ser elaborado em língua inglesa, seguindo tanto quanto possível a redação utilizada no Resumo. \lipsum[2]
}

\keywords{words. with. period. as separator.}


\begin{document}
\section{Introduction}


\lipsum[2]

\section{Section}

\lipsum[3]

%\begin{figure}[h]\centering
%\includegraphics[width=4cm]
%{example-image-a}
%\caption{Legenda aqui}
%\end{figure}

\begin{figure}[h]\centering
\includegraphics[width=8cm]
{logo_rainbow}
\caption{Legenda aqui}
\end{figure}

\section{Section}
\begin{equation}
    \frac{y}{x}=2
\end{equation}
\begin{align}
    xyz=\delta
\end{align}
\lipsum[6]

\subsection{Subsection}

\lipsum[6]




%\Referenciasbibliograficas

Aqui citamos como exemplo \cite{Boldrini}, \cite{Cuminato} e \cite{daSilva}

\begin{thebibliography}{100}

\bibitem{Boldrini} 
Boldrini, J. L., Costa, S. I. R., Ribeiro, V. R. and Wetzler, H. G. {\it Álgebra Linear e Aplicações, 3a. edição}. Harbra, São Paulo, 1984.

\bibitem{Cuminato}
Cuminato, J. A. and Ruas, V. Unification of distance inequalities for linear variational problems, 
{\it Comp. Appl. Math.}, 2014. DOI: 10.1007/s40314-014-0163-6.
 
\bibitem{daSilva} 
Da Silva, P. L. and Freire, I. L. On the group analysis of a modified Novikov equation, 
{\it Interdisciplinary Topics in Applied Mathematics, Modeling and Computational Science}, 
Springer Proceedings in Mathematics and Statistics,  volume 117, chapter 23, pages 161-166, 2015.

\bibitem{Diniz1}
Diniz, G. L. A mudança no habitat de populações de peixes: de rio a represa -- o modelo 
matemático, Dissertação de Mestrado, Unicamp, 1994.

\bibitem{Diniz2}
Diniz, G. L., Meyer, J. F. C. A. e Barros, L. C. Solução numérica para um problema de 
Cauchy Fuzzy que modela o decaimento radioativo, {\it TEMA},  23:63--72, 2001. DOI:10.1007/s40314-014-0163-6.

\bibitem{Gomes}
Gomes, L. T., De Barros, L. C. and Bede, B. Fuzzy differential equation in various approaches. 
In {\it SpringerBriefs in Mathematics}. SBMAC- Springer, 2015. ISSN: 2191-8198.

%\bibitem{Jafelice} Jafelice, R. M., Barros, L. C. and Bassanezi, R. C. Study of the dynamics of HIV under treatment considering fuzzy delay, {\it Comp. Appl. Math.}, 33:45--61, 2014.

%\bibitem{Mallet}
%Mallet, S. M. Análise Numérica de Elementos Finitos. Tese de Doutorado, LNCC/MCTI, 1990.

\bibitem{Santos} Santos, I. L. D. e Silva, G. N. Uma classe de problemas de controle ótimo 
em escalas temporais, {\it Proceeding Series of the Brazilian Society of Computational and 
Applied Mathematics}, volume 1, 2013. DOI: 10.5540/03.2013.001.01.0177.

\end{thebibliography}

%Devem ser alinhadas a esquerda, espaçamento simples e separadas entre si por uma linha em branco. Utilizar o recurso negrito para destacar o elemento título de cada publicação referenciada.

%Veja exemplos-base no arquivo ``Instruções Específicas sobre a submissão e formatação do artigo'' (3.6).


\end{document}